% GENERATED FILE. Edit canonical lessons, then run npm run generate:books.
% canonical-source-hash: fnv1a64:865f61fc902c025a
% canonical-lessons: JA-C183-kyaku, JA-C183-tenin, JA-C183-uriba, JA-C183-erabu, JA-C183-narabu, JA-R183-review-178-183

\chapter{Customers and Shop Staff}
\label{ch:ja-customers-and-shop-staff}
\hlchaptermodality{183}

\begin{chapteropening}
\textbf{By the end of this chapter:} \emph{I can talk about customers, staff and choosing goods.}

Complete the last lesson of chapter 183: I can talk about customers, staff and choosing goods.
\end{chapteropening}

\section[\texorpdfstring{\ja{きゃく} (kyaku)}{きゃく (kyaku)}]{\ja{きゃく} (kyaku) — customer; guest, visitor}

\label{lesson:JA-C183-kyaku}



\begin{quote}
Before the new one: say the Japanese for banknote, then the Japanese for small coins.
\end{quote}

\subsection*{You'll want to know: \ja{きゃく}}
\textbf{\ja{きゃく}} — \emph{kyaku} — \textquotedblleft{}customer; guest, visitor\textquotedblright{}.

\begin{quote}
\textbf{\ja{この} \ja{みせは} \ja{きゃくが} \ja{おおいです。}} — \emph{kono mise wa kyaku ga ōi desu} — \textquotedblleft{}This shop has many customers.\textquotedblright{}
\end{quote}

\begin{cousinweb}
\textbf{\ja{きゃく}} (\emph{kyaku}) is Sino-Japanese (kango): one Chinese character, \textquotedblleft{}guest, customer\textquotedblright{}, read with its on'yomi, a reading borrowed from Chinese (Middle Chinese \emph{khaek}) (KANJIDIC2; Middle Chinese in Baxter's transcription).
\end{cousinweb}

\begin{grammarlens}[title={What you've built}]
One more word for everyday life at A2.
\end{grammarlens}

\subsection*{Your turn}
\begin{itemize}
\raggedright
  \item \emph{Say it:} \emph{kyaku}
  \item \emph{Say it:} \emph{kyaku}, once more
  \item \emph{Say it:} \emph{kono mise wa kyaku ga ōi desu} — This shop has many customers.
  \item \emph{Recall:} say the Japanese for banknote, then the Japanese for small coins, then say \emph{kyaku} again
\end{itemize}

\begin{tcolorbox}[breakable,colback=teal!4,colframe=teal!35!black,title={Before you move on}]
What does \ja{きゃく} mean? (\textquotedblleft{}Customer; guest, visitor\textquotedblright{}.) Say it once more.
\end{tcolorbox}
\section[\texorpdfstring{\ja{てんいん} (ten'in)}{てんいん (ten'in)}]{\ja{てんいん} (ten'in) — shop assistant; clerk}

\label{lesson:JA-C183-tenin}



\begin{quote}
Before the new one: say the Japanese for small coins, then the Japanese for customer.
\end{quote}

\subsection*{You'll want to know: \ja{てんいん}}
\textbf{\ja{てんいん}} — \emph{ten'in} — \textquotedblleft{}shop assistant; clerk\textquotedblright{}.

\begin{quote}
\textbf{\ja{あの} \ja{てんいんは} \ja{しんせつです。}} — \emph{ano ten'in wa shinsetsu desu} — \textquotedblleft{}That shop assistant is kind.\textquotedblright{}
\end{quote}

\begin{cousinweb}
\textbf{\ja{てんいん}} (\emph{ten'in}) is Sino-Japanese (kango): two Chinese characters read with their on'yomi, readings borrowed from Chinese. They are \emph{ten}, \textquotedblleft{}shop\textquotedblright{} (Middle Chinese \emph{temH}) and \emph{in}, \textquotedblleft{}member, staff\textquotedblright{} (\emph{hjwen}) (KANJIDIC2; Middle Chinese in Baxter's transcription).
\end{cousinweb}

\begin{grammarlens}[title={What you've built}]
One more word for everyday life at A2.
\end{grammarlens}

\subsection*{Your turn}
\begin{itemize}
\raggedright
  \item \emph{Say it:} \emph{ten'in}
  \item \emph{Say it:} \emph{ten'in}, once more
  \item \emph{Say it:} \emph{ano ten'in wa shinsetsu desu} — That shop assistant is kind.
  \item \emph{Recall:} say the Japanese for small coins, then the Japanese for customer, then say \emph{ten'in} again
\end{itemize}

\begin{tcolorbox}[breakable,colback=teal!4,colframe=teal!35!black,title={Before you move on}]
What does \ja{てんいん} mean? (\textquotedblleft{}Shop assistant; clerk\textquotedblright{}.) Say it once more.
\end{tcolorbox}
\section[\texorpdfstring{\ja{うりば} (uriba)}{うりば (uriba)}]{\ja{うりば} (uriba) — sales floor; counter, department}

\label{lesson:JA-C183-uriba}



\begin{quote}
Before the new one: say the Japanese for customer, then the Japanese for shop assistant.
\end{quote}

\subsection*{You'll want to know: \ja{うりば}}
\textbf{\ja{うりば}} — \emph{uriba} — \textquotedblleft{}sales floor; counter, department\textquotedblright{}.

\begin{quote}
\textbf{\ja{くつの} \ja{うりばは} \ja{どこですか。}} — \emph{kutsu no uriba wa doko desu ka} — \textquotedblleft{}Where is the shoe department?\textquotedblright{}
\end{quote}

\begin{cousinweb}
\textbf{\ja{うりば}} (\emph{uriba}) is native Japanese, built from \emph{uri}, the stem of \emph{uru}, \textquotedblleft{}to sell\textquotedblright{}, and \emph{ba}, \textquotedblleft{}place\textquotedblright{} (JMdict, KANJIDIC2).
\end{cousinweb}

\begin{grammarlens}[title={What you've built}]
One more word for everyday life at A2.
\end{grammarlens}

\subsection*{Your turn}
\begin{itemize}
\raggedright
  \item \emph{Say it:} \emph{uriba}
  \item \emph{Say it:} \emph{uriba}, once more
  \item \emph{Say it:} \emph{kutsu no uriba wa doko desu ka} — Where is the shoe department?
  \item \emph{Recall:} say the Japanese for customer, then the Japanese for shop assistant, then say \emph{uriba} again
\end{itemize}

\begin{tcolorbox}[breakable,colback=teal!4,colframe=teal!35!black,title={Before you move on}]
What does \ja{うりば} mean? (\textquotedblleft{}Sales floor; counter, department\textquotedblright{}.) Say it once more.
\end{tcolorbox}
\section[\texorpdfstring{\ja{えらぶ} (erabu)}{えらぶ (erabu)}]{\ja{えらぶ} (erabu) — to choose; to pick}

\label{lesson:JA-C183-erabu}



\begin{quote}
Before the new one: say the Japanese for shop assistant, then the Japanese for sales floor.
\end{quote}

\subsection*{You'll want to know: \ja{えらぶ}}
\textbf{\ja{えらぶ}} — \emph{erabu} — \textquotedblleft{}to choose; to pick\textquotedblright{}.

\begin{quote}
\textbf{\ja{あには} \ja{いつも} \ja{やすい} \ja{くつを} \ja{えらびます。}} — \emph{ani wa itsumo yasui kutsu o erabimasu} — \textquotedblleft{}My older brother always chooses cheap shoes.\textquotedblright{}
\end{quote}

Polite \textbf{\ja{えらびます}}, \emph{erabimasu}; a godan verb.

\begin{cousinweb}
\textbf{\ja{えらぶ}} (\emph{erabu}) is a native Japanese word (wago), inherited rather than borrowed from Chinese; its kanji was attached later as a kun'yomi (JMdict, KANJIDIC2).
\end{cousinweb}

\begin{grammarlens}[title={What you've built}]
One more word for everyday life at A2.
\end{grammarlens}

\subsection*{Your turn}
\begin{itemize}
\raggedright
  \item \emph{Say it:} \emph{erabu}
  \item \emph{Say it:} \emph{erabu}, once more
  \item \emph{Say it:} \emph{ani wa itsumo yasui kutsu o erabimasu} — My older brother always chooses cheap shoes.
  \item \emph{Recall:} say the Japanese for shop assistant, then the Japanese for sales floor, then say \emph{erabu} again
\end{itemize}

\begin{tcolorbox}[breakable,colback=teal!4,colframe=teal!35!black,title={Before you move on}]
What does \ja{えらぶ} mean? (\textquotedblleft{}To choose; to pick\textquotedblright{}.) Say it once more.
\end{tcolorbox}
\section[\texorpdfstring{\ja{ならぶ} (narabu)}{ならぶ (narabu)}]{\ja{ならぶ} (narabu) — to line up; to queue}

\label{lesson:JA-C183-narabu}



\begin{quote}
Before the new one: say the Japanese for sales floor, then the Japanese for to choose.
\end{quote}

\subsection*{You'll want to know: \ja{ならぶ}}
\textbf{\ja{ならぶ}} — \emph{narabu} — \textquotedblleft{}to line up; to queue\textquotedblright{}.

\begin{quote}
\textbf{\ja{きゃくが} \ja{みせの} \ja{まえに} \ja{ならびます。}} — \emph{kyaku ga mise no mae ni narabimasu} — \textquotedblleft{}Customers line up in front of the shop.\textquotedblright{}
\end{quote}

Polite \textbf{\ja{ならびます}}, \emph{narabimasu}; a godan verb.

\begin{cousinweb}
\textbf{\ja{ならぶ}} (\emph{narabu}) is a native Japanese word (wago), inherited rather than borrowed from Chinese; its kanji was attached later as a kun'yomi (JMdict, KANJIDIC2).
\end{cousinweb}

\begin{grammarlens}[title={What you've built}]
One more word for everyday life at A2.
\end{grammarlens}

\subsection*{Your turn}
\begin{itemize}
\raggedright
  \item \emph{Say it:} \emph{narabu}
  \item \emph{Say it:} \emph{narabu}, once more
  \item \emph{Say it:} \emph{kyaku ga mise no mae ni narabimasu} — Customers line up in front of the shop.
  \item \emph{Recall:} say the Japanese for sales floor, then the Japanese for to choose, then say \emph{narabu} again
\end{itemize}

\begin{tcolorbox}[breakable,colback=teal!4,colframe=teal!35!black,title={Before you move on}]
What does \ja{ならぶ} mean? (\textquotedblleft{}To line up; to queue\textquotedblright{}.) Say it once more.
\end{tcolorbox}
\section[(dialogue)]{Words from chapters 178–183 — a review}

\label{lesson:JA-R183-review-178-183}



\begin{quote}
Say the Japanese for pork and beef.
\end{quote}

\subsection*{Your turn}
Say these aloud:

\begin{itemize}
\raggedright
  \item pork, beef, chicken meat, milk, to chill
  \item black tea, rice wine, to pour, to spill, to warm up
  \item cafe, cafeteria, miso soup, to grip, rice ball
  \item shopping, price, how much, to sell, to sell well
  \item to pay, change, cash, banknote, small coins
  \item customer, shop assistant, sales floor, to choose, to line up
\end{itemize}

\begin{tcolorbox}[breakable,colback=teal!4,colframe=teal!35!black,title={Before you move on}]
Pork? (\textbf{\ja{ぶたにく}}, \emph{butaniku}.) Beef? (\textbf{\ja{ぎゅうにく}}, \emph{gyūniku}.) Chicken meat? (\textbf{\ja{とりにく}}, \emph{toriniku}.) Milk? (\textbf{\ja{ぎゅうにゅう}}, \emph{gyūnyū}.) To chill? (\textbf{\ja{ひやす}}, \emph{hiyasu}.) Black tea? (\textbf{\ja{こうちゃ}}, \emph{kōcha}.) Rice wine? (\textbf{\ja{おさけ}}, \emph{osake}.) To pour? (\textbf{\ja{そそぐ}}, \emph{sosogu}.) To spill? (\textbf{\ja{こぼす}}, \emph{kobosu}.) To warm up? (\textbf{\ja{あたためる}}, \emph{atatameru}.) Cafe? (\textbf{\ja{きっさてん}}, \emph{kissaten}.) Cafeteria? (\textbf{\ja{しょくどう}}, \emph{shokudō}.) Miso soup? (\textbf{\ja{みそしる}}, \emph{misoshiru}.) To grip? (\textbf{\ja{にぎる}}, \emph{nigiru}.) Rice ball? (\textbf{\ja{おにぎり}}, \emph{onigiri}.) Shopping? (\textbf{\ja{かいもの}}, \emph{kaimono}.) Price? (\textbf{\ja{ねだん}}, \emph{nedan}.) How much? (\textbf{\ja{いくら}}, \emph{ikura}.) To sell? (\textbf{\ja{うる}}, \emph{uru}.) To sell well? (\textbf{\ja{うれる}}, \emph{ureru}.) To pay? (\textbf{\ja{はらう}}, \emph{harau}.) Change? (\textbf{\ja{おつり}}, \emph{otsuri}.) Cash? (\textbf{\ja{げんきん}}, \emph{genkin}.) Banknote? (\textbf{\ja{さつ}}, \emph{satsu}.) Small coins? (\textbf{\ja{こぜに}}, \emph{kozeni}.) Customer? (\textbf{\ja{きゃく}}, \emph{kyaku}.) Shop assistant? (\textbf{\ja{てんいん}}, \emph{ten'in}.) Sales floor? (\textbf{\ja{うりば}}, \emph{uriba}.) To choose? (\textbf{\ja{えらぶ}}, \emph{erabu}.) To line up? (\textbf{\ja{ならぶ}}, \emph{narabu}.)
\end{tcolorbox}
